\documentclass[11pt,a4paper]{article}
\usepackage[T1]{fontenc}
\usepackage[utf8]{inputenc}
\usepackage{lmodern,amsmath,amssymb,amsthm,mathtools}
\usepackage[margin=25mm]{geometry}
\usepackage{booktabs,longtable,array,microtype,xcolor,enumitem,listings,hyperref}
\hypersetup{colorlinks=true,linkcolor=blue!45!black,urlcolor=blue!45!black,citecolor=blue!45!black,
 pdftitle={Smaller Diophantine forms of ZFC consistency},pdfauthor={Research contribution prepared for ProveIt}}
\setlength{\emergencystretch}{3em}
\setlength{\parskip}{3pt}
\setlist{itemsep=2pt,topsep=4pt}
\newtheorem{theorem}{Theorem}[section]
\newtheorem{lemma}[theorem]{Lemma}
\newtheorem{proposition}[theorem]{Proposition}
\newtheorem{corollary}[theorem]{Corollary}
\theoremstyle{definition}\newtheorem{definition}[theorem]{Definition}
\theoremstyle{remark}\newtheorem{remark}[theorem]{Remark}
\newcommand{\Z}{\mathbb Z}
\newcommand{\N}{\mathbb N}
\newcommand{\Pos}{\mathbb Z_{>0}}
\newcommand{\Con}{\operatorname{Con}}
\newcommand{\pc}{\operatorname{pc}}
\newcommand{\vv}{\nu_2}
\newcommand{\code}[1]{\texttt{\detokenize{#1}}}
\newcommand{\pin}{\texttt{3b201cdea32d2c816f5f3d20734051f25fd3ef14}}
\newcommand{\MF}{M_F^{\mathrm{src}}}
\lstset{basicstyle=\ttfamily\footnotesize,breaklines=true,columns=fullflexible,keepspaces=true,frame=single}
\title{\textbf{Smaller Diophantine forms of ZFC consistency}\\[5pt]
 \large Combining blank-input compilation, ordinary Pell norms,\\
 root elimination, and a five-cell local language}
\author{Research contribution prepared for Vladimir Reshetnikov's ProveIt repository}
\date{22 October 2026\\[4pt]\small Audited repository snapshot: \pin}
\begin{document}
\maketitle
\begin{abstract}
We combine and simplify the eight fixed-instance constructions collected in
ProveIt's \emph{Con(ZFC) in Diophantine form} report. The principal transparent
construction is a system of six polynomial equations in fifteen positive
unknowns, of maximum degree 28, with a 60-operation arithmetic evaluator.
An auxiliary-gap substitution lowers the maximum degree to 22 at a cost
of one additional operation; its sum of squares has degree 44 and costs
78 operations. Combining this gap with Pell-root elimination gives one
degree-219 equation in fourteen positive unknowns at 75 operations.
A 73-operation, degree-342 alternative offers another tradeoff.
A quadratic presentation uses 32 positive unknowns and 23 equations,
yielding one quartic.

A separate simplification removes the four unused corners of the local
computation windows. The selector alphabet drops from approximately
\(1.42\cdot10^{21}\) to exactly \(38\,256\,607\,897\), and the bit length
of the base coefficient drops from \(5^{67}\) to \(5^{36}\). All four
coefficients are fixed by a terminating finite recipe from the literal
K215 transition table. They are not free parameters and are not expanded.

For each presented system or polynomial, ordinary mathematics formalizable
in ZFC proves that it has a solution in its stated domain if and only if
ZFC has a finite refutation. Thus ZFC proves its unsolvability equivalent
to \(\Con(\mathrm{ZFC})\). The proof includes a sign-independent auxiliary
Pell argument and a main-root sign argument; these reduce several
unrestricted-integer conversions substantially. We give exact circuits,
quadratic residuals, source hashes, reproducible finite audits, and an
explicit account of the inherited proof dependencies. No global
optimality, consistency proof, or new Lean/Rocq verification is claimed.
\end{abstract}
\clearpage
\begingroup\small\setlength{\parskip}{0pt}\tableofcontents\endgroup
\clearpage

\section{What is improved, and what ``simpler'' means}
\label{sec:overview}

The requested object is a fixed polynomial or finite system, with integer
coefficients, for which
\begin{equation}\label{eq:target}
 \mathrm{ZFC}\vdash
 \bigl[\text{no solution in the stated domain}\bigr]
 \ \longleftrightarrow\ \Con(\mathrm{ZFC}).
\end{equation}
There are several incompatible notions of simplicity: number of unknowns,
number of equations, degree, arithmetic evaluation length, and coefficient
height. We give a small family of concrete improvements and do not identify
one of these objectives with all the others.

An arithmetic operation means one binary addition, subtraction, or
multiplication, including an operation with a fixed numeral. Reading a
port, reusing a computed value, and comparing two values in a system
cost zero. Every source packet gives the complete acyclic circuit under
this convention. Coefficients are fixed numerals; the table does not
charge their bit lengths. The separate coefficient analysis does.

The following primary presentations use strictly positive unknowns.
Their coefficients are the \emph{same four new integers} from
Section~\ref{sec:coefficients}; using the old package-06 coefficients
instead also preserves the arithmetic conclusions.
\begin{center}\small
\begin{tabular}{@{}lrrrr@{}}
\toprule
Presentation & Unknowns & Equations & Degree & Operations\\
\midrule
Gap root norm \(\Phi_{14}^{\rm gap}\)&14&1&219&75\\
Root outer norm \(\Psi_{14}\)&14&1&342&73\\
Ordinary system \(\mathcal S_{60}\) &15&6&28 maximum&60\\
Gap system &15&6&22 maximum&61\\
Product \(P_{69}\) &15&1&99&69\\
Gap product &15&1&93&70\\
Sum of squares \(S_{77}\) &15&1&56&77\\
Gap sum of squares &15&1&44&78\\
Quadratic system \(\mathcal Q_{32}\) &32&23&2 maximum&not optimized\\
Sum of squares of \(\mathcal Q_{32}\) &32&1&4&not optimized\\
\bottomrule
\end{tabular}
\end{center}
The fourteen-variable equations are displayed and measured in
Sections~\ref{sec:fourteen} and~\ref{sec:gap14}. The integer-domain
presentations appear in Section~\ref{sec:integer}. In particular, there
is a polynomial in 32 unrestricted integer unknowns of degree 434,
one in 33 of degree 88, and one in 50 of degree 8.

\subsection{Comparison of the eight source packages}

The repository reconciliation records eight independently prepared
packages, all based on the earlier snapshot
\code{cdd582bec13f908df3711fa9e05d1ddac368a5db}. We inspected the
reconciled snapshot \pin. The following table selects the source
constructions that affect the present comparison; it is not a claim that
every historical check or formal build was rerun here. Full provenance
and the source paths are recorded in the archive and in the reconciliation
note~\cite{collection}.

\begin{center}\small
\begin{tabular}{@{}>{\raggedright\arraybackslash}p{28mm}>{\raggedright\arraybackslash}p{83mm}>{\raggedright\arraybackslash}p{30mm}@{}}
\toprule
Source & Relevant construction & Role here\\
\midrule
01-explicit82 &18 variables, one degree-187 equation, 82 operations;
 fixed-input specialization of the full compiler &Baseline specialization\\
02-system63 &16 variables, seven equations, maximum degree 28,
 63 operations; ordinary strong norm &Low-degree Pell block\\
03-system64 &A closely related 64-operation system; elimination
 of the strong root gives 15 variables &Root-elimination method\\
04-poly73 &16 variables, degree 156, 73 operations; direct
 two-sided normalizer and affine input guard &Alternative normalization\\
05-fixed15-16 &15-variable norm products: 82 operations/degree 295
 and 80 operations/degree 346 &Two root finalizers\\
06-blank67 &Raw K215; 15 variables, six equations/61 operations;
 polynomial 67 operations/degree 155 &Blank semantics and source circuit\\
07-fixed70 &16-variable product, 70 operations/degree 169;
 17-variable system, 61 operations/degree 28 &Proof-order and scale checks\\
08-start-only67 &Same 67-operation arithmetic skeleton with different
 coefficients; quartic in 34 positive variables &Quadratic lifting\\
\bottomrule
\end{tabular}
\end{center}

The gains combine compatible interfaces, rather than numerical
coefficients from different machines. Package 02's ordinary strong norm
can be used after package 06 deletes the entire input interface. Package
08's lifting then loses two registers. Packages 03 and 05's root
elimination can be applied \emph{after} deletion of the input interface,
giving fourteen unknowns. Finally, the local predicate in package 06
does not inspect its four corner entries. Removing those entries
reduces coefficient height without adding any arithmetic gate.
The additional substitution \(y=V+e\) cancels the highest auxiliary
square terms. Applying that substitution before root elimination gives
the degree-219 fourteen-variable equation, at the same 75 operations
as its degree-293 predecessor.

The 60-operation system improves the earlier 16-variable, seven-equation,
63-operation system at the same maximum degree 28. It also improves the
61-operation, 15-variable blank system in both operations and degree.
The quartic improves 34 positive unknowns and 25 quadratic equations to
32 and 23. The 67-operation source remains useful when operation count
alone is prioritized; our 69-operation product trades two operations for
a degree decrease from 155 to 99.

Classical low-variable universal constructions are a different comparison.
Jones's work includes nine-variable universal representations with a very
different degree and coefficient tradeoff~\cite{jones}. The repository's
\code{Jones1982.nine_unknowns} states an existence theorem, not a literal
nine-variable specialization of these four ports. Modern relation-combining
constructions also give representations in eleven unrestricted integer
unknowns, with much larger degrees~\cite{bayer2025}. We make no absolute
minimum-variable or discovery-priority claim.

\section{The six equations in fifteen unknowns}
\label{sec:system}

Fix the four integers
\begin{equation}\label{eq:ports}
 \beta=B-1,\qquad K=D_C+BD_R,\qquad M_C,\qquad
 \MF=M_F+B-1
\end{equation}
from Section~\ref{sec:coefficients}. Here \(M_F\) is the native mask;
\(\MF\) is the shifted source port. They must not be confused.

The fifteen supplied unknowns are
\[
 (J,F,\alpha,z,f,h,i,T,s,w,\tau,\eta,\zeta,y,\gamma)\in\Pos^{15}.
\]
Define the following \emph{polynomial abbreviations}. They are not
additional existential unknowns:
\begin{align}
 q&=\beta J+1,& C&=q-F-\alpha,&X&=wq,&Y&=sq^3,\notag\\
 E&=XY,&k&=\eta+\zeta,&c&=kY+\eta,&a&=E+Y,\notag\\
 H&=4a+3,&\Delta&=a^2+H=(a+2)^2-1,&D&=X+ac+\gamma H.
 \label{eq:definitions}
\end{align}
Further set
\begin{align}
 G&=(q-1)(q-F)+\alpha=q(q-F)-C,\notag\\
 R&=G(q^2-1)+(M_C+q\MF)J,\notag\\
 L&=EkY,\qquad V=c(Tf-1)-Rf^2,\qquad
 I=ic^2,\qquad U=\Delta(f^2-1).\label{eq:packing}
\end{align}
The system is
\begin{equation}\boxed{\begin{aligned}
 \tau^2-L(L+k)&=1,\\
 D^2-\Delta c^2&=1,\\
 U(V^2-y^2)+y^2&=1,\\
 k-hE-R&=1,\\
 (K+w)C&=F+z(q-1),\\
 I^2&=U.
\end{aligned}}\label{eq:system60}
\end{equation}
All powers have fixed integer exponents. Evaluation uses only ordinary
integer polynomial arithmetic. The source
\code{arithmetic/ordinary_blank_system60.json} has 60 live gates;
all residuals have degree at most 28.

\begin{theorem}\label{thm:main}
For the four coefficients in Section~\ref{sec:coefficients}, ZFC proves
\[
 \exists\mathbf x\in\Pos^{15}\;\mathcal S_{60}(\mathbf x)
 \quad\longleftrightarrow\quad
 \exists p\;\operatorname{Proof}_{\mathrm{ZFC}}(p,\bot).
\]
Thus the absence of a positive solution is provably equivalent in ZFC to
\(\Con(\mathrm{ZFC})\). The same assertion holds for the other
presentations in this article, in their explicitly stated domains.
\end{theorem}

\subsection{Two single-polynomial presentations}

For the product presentation use the same definitions but call the
transport unknown \(z_0\), and put
\begin{align*}
 N_0&=\tau^2-L(L+k),&N_1&=D^2-\Delta c^2,\\
 N_A&=U(V^2-y^2)+y^2,&N_k&=k-hE-R,\\
 N_t&=(K+w)C+q-F-z_0(q-1),&N_s&=1+I^2-U.
\end{align*}
Then
\begin{equation}\boxed{P_{69}=N_0N_1N_AN_kN_tN_s-1=0.}
\label{eq:p69}\end{equation}
Its complete evaluator has 69 operations and exact degree 99.
The product and the six-equation system correspond by
\(z_0=z+1\). The inverse shift preserves positivity on solutions;
this is proved below and is not assumed on arbitrary tuples.

Alternatively, let \(r_1,\ldots,r_6\) be the six left-minus-right
residuals in \eqref{eq:system60}. Then
\begin{equation}\boxed{S_{77}=r_1^2+\cdots+r_6^2=0}
\label{eq:sos77}\end{equation}
has exactly the same zero set over the integers in the stipulated domain.
The six residual subtractions, six squarings and five additions add 17
operations to the 60-operation comparison circuit. Its exact degree is 56.

\section{Pell facts and the signed auxiliary argument}
\label{sec:pell}

For \(A\ge2\) define
\[
 \chi_A(n)+\psi_A(n)\sqrt{A^2-1}
  =(A+\sqrt{A^2-1})^n\quad(n\ge0).
\]
Both sequences have integer coefficients and obey
\(u_{n+2}=2Au_{n+1}-u_n\). Their initial pairs are
\((\chi_0,\chi_1)=(1,A)\) and
\((\psi_0,\psi_1)=(0,1)\). Positive norm-one solutions are exactly
these pairs. The usual proof divides by the least positive unit until
the quotient lies between 1 and that unit. This takes place in the
integer-coefficient quadratic ring and needs no assertion about the
full ring of algebraic integers.

\subsection{The ordinary strong norm really does recover integrality}

The equation \((ic^2)^2=\Delta(f^2-1)\) cannot simply be read as an
integral Pell equation at \(A\): divisibility by \(\Delta\) must be
proved. We use the following precise inherited lemma~\cite{rank}, giving
its argument to make the dependency explicit.

\begin{lemma}[Ordinary strong rank]\label{lem:rank}
Let \(\Delta=A^2-1\), \(c=\psi_A(p)>A\Delta^2\), \(p\ge1\),
and \(i>0\). If \((ic^2)^2=\Delta(f^2-1)\) for an integer \(f\),
then for an integer \(m>0\),
\[
 |f|=\chi_A(m),\quad p\mid m,\quad c\mid m,\quad
 ic^2=\Delta\psi_A(m),\quad c^2\mid f^2-1.
\]
\end{lemma}
\begin{proof}
Write \(\Delta=d_0s_0^2\) with \(d_0>1\) squarefree. Let
\(F+y_0\sqrt{d_0}\) be the least positive integer-coefficient unit
of norm one. For some \(e\ge1\),
\[
 A=\chi_F(e),\quad s_0=y_0S,\quad S=\psi_F(e),\quad
 \Delta=(F^2-1)S^2,\quad \psi_F(ep)=Sc.
\]
Put \(W=ic^2\). Divisibility in \(W^2=\Delta(f^2-1)\) first
gives \(s_0\mid W\), then \(d_0\mid W/s_0\). Thus
\(|f|+(W/(d_0s_0))\sqrt{d_0}\) is an integral unit, so for some
\(n>0\),
\[
 |f|=\chi_F(n),\qquad W=(F^2-1)S\psi_F(n).
\]
The Pell sequence is strongly divisible:
\(\gcd(\psi_F(u),\psi_F(v))=\psi_F(\gcd(u,v))\).
This follows from the addition identity, coprimality of the two
coordinates of a norm-one unit, and the Euclidean algorithm.

Let \(g=\gcd(n,ep)\), \(t=\psi_F(g)\), and
\(T_0=(F^2-1)S\le\Delta\). Since \(c^2\mid T_0\psi_F(n)\),
\[
 c\le\gcd(c^2,\Delta c)
   \mid\gcd(T_0\psi_F(n),T_0\psi_F(ep))=T_0t,
\]
so \(t\ge c/\Delta\). If \(g<ep\), then \(2g\le ep\).
The duplication identity gives
\(2t^2<\psi_F(2g)\le Sc<Ac\), since \(S<A\).
But \(t^2\ge c^2/\Delta^2>Ac\), a contradiction. Therefore
\(ep\mid n\). Writing \(n=em\) gives \(p\mid m\) and
\(W=\Delta\psi_A(m)\).

Write \(m=pu\). Expansion of
\((\chi_A(p)+c\sqrt\Delta)^u\) modulo \(c^2\) gives
\[
 \psi_A(pu)/c\equiv u\chi_A(p)^{u-1}\pmod c.
\]
Thus \(c\mid\Delta u\). Also
\(\gcd(c,\Delta)=\gcd(p,\Delta)\), by the Pell recurrence
modulo \(\Delta\). This gcd divides \(p\), whence \(c\mid pu=m\).
Finally \(f^2-1=\Delta\psi_A(m)^2=ic^2\psi_A(m)\), proving
the last divisibility.
\end{proof}

\subsection{A step-down lemma with all signs visible}

\begin{lemma}[Strict chi step-down]\label{lem:stepdown}
If \(A\ge2\), \(n\ge0\), \(0<r<m\), and
\(\chi_A(n)\equiv\chi_A(r)\pmod{\chi_A(m)}\), then
\(n\equiv r\) or \(-r\pmod{2m}\).
\end{lemma}
\begin{proof}
Modulo \(\chi_A(m)\), the addition identities give
\(\chi_A(n+2m)\equiv-\chi_A(n)\) and
\(\chi_A(2m-n)\equiv-\chi_A(n)\). Reducing and reflecting the
index gives a residue of the form \(\pm\chi_A(t)\) with
\(0\le t\le m\). The case \(t=m\) is zero and cannot work.
For \(t,r<m\), both values are positive and less than
\(\chi_A(m)/2\), since \(\chi_A(m)>2\chi_A(m-1)\).
Their sum is smaller than the modulus, so the negative sign is
impossible. Equality of the positive residues gives \(t=r\) by
strict growth. Undoing the reductions proves the assertion.
\end{proof}

\begin{lemma}[Signed auxiliary index]\label{lem:signed}
Suppose \(c=\psi_A(p)\), \(|f|=\chi_A(m)\), \(c\mid m\),
\(m>2p\), and \(0<p,R<c/2\). Let \(I>1\) be divisible by \(c\),
with \(I^2=\Delta(f^2-1)\). If integers \(V,y\) satisfy
\[
 I^2V^2-(I^2-1)y^2=1,\quad
 V\equiv-c\pmod{|f|},\quad V\equiv-R\pmod c,
\]
then \(p=R\). No sign condition on \(V\) or \(y\) is needed.
\end{lemma}
\begin{proof}
Neither \(V\) nor \(y\) can vanish. Pell classification gives
\(I|V|=\chi_I(\ell)\), \(|y|=\psi_I(\ell)\).
Modulo \(I\), even-index chi values are \(\pm1\), so
\(\ell=2v+1\) is odd. There is an integer polynomial \(Q_v\) with
\[
 \chi_t(2v+1)=tQ_v(t^2),\quad
 Q_v(1-A^2)=(-1)^v\psi_A(2v+1),\quad
 Q_v(0)=(-1)^v(2v+1).
\]
These identities follow from the second-order recurrence (equivalently,
the odd Chebyshev identity). Let \(\sigma=\operatorname{sign}V\);
then \(|V|=\sigma V=Q_v(I^2)\). As \(I^2\equiv1-A^2\pmod{|f|}\),
the first congruence gives \(\psi_A(\ell)^2\equiv c^2\pmod{|f|}\).
Using \(\chi_A(2t)=1+2\Delta\psi_A(t)^2\), Lemma~\ref{lem:stepdown}
gives \(\ell\equiv\pm p\pmod m\), hence modulo \(c\).
Modulo \(c\), \(I^2\equiv0\) and the second congruence gives
\(\sigma(-1)^v\ell\equiv-R\pmod c\). Thus
\(R\equiv\pm p\pmod c\). The strict bounds exclude both a nonzero
wrap and the negative alternative, leaving \(R=p\).
\end{proof}

The lemma explains why one should not require the deleted auxiliary
quotients to be positive before extracting the index. It also prevents
an invalid use of a positive-root parity lemma after signs have been
released. We recover oddness from the first/main index relation; any
stronger residue class is recovered later from the actual masks.

\section{Proof of the arithmetic reduction}
\label{sec:arithmetic-proof}

We prove soundness on the larger domain in which only
\begin{equation}\label{eq:guarded}
 J,F,\alpha,z,i,s,w,\eta,\zeta>0
\end{equation}
are required positive. The six remaining coordinates
\(h,f,T,\tau,y,\gamma\) may be arbitrary integers. For the product,
replace \(z\) by the positive \(z_0\). This stronger theorem supplies
the integer conversions without inserting unnecessary positivity guards.

\subsection{Units, the boundary case, and pre-decoding bounds}

For the system, transport immediately implies \(C>0\). For the product,
each of its six integer factors is \(\pm1\). The first factor cannot
be \(-1\): the negative Pell equation for \(u(u+1)\), \(u=XY^2>1\),
has no solution. A least positive second coordinate would be decreased
by multiplication by \((2u+1)-2\sqrt{u(u+1)}\).

Since \(\Delta=(a+2)^2-1\) is 0 or 3 modulo 4, the main norm cannot
be \(-1\). Nor can \(N_s=1+I^2-\Delta(f^2-1)\): modulo 4 it is
either \(1+I^2\) or \(I^2+f^2\), neither of which is 3.
Thus \(N_s=1\), giving \(U=I^2>0\). The auxiliary factor is now
congruent to either \(V^2\) or \(y^2\) modulo 4, and also equals 1.
Only \(N_k=\varepsilon\) and \(N_t=\nu\) remain, with
\(\varepsilon,\nu\in\{1,-1\}\) and \(\varepsilon\nu=1\).

Weak transport gives \(C\ge0\), since
\[
 q-F-z_0(q-1)=1-F-(z_0-1)(q-1)\le0
\]
and \(C\le-1\) would force \(N_t<-1\). We retain \(C=0\)
at this stage. It really is an algebraically possible boundary for the
negative transport unit and must not be discarded by assertion.

In both presentations, \(q\ge B\ge16\), \(F\ge1\), \(\alpha\ge1\),
and \(q-F=C+\alpha>0\). Therefore
\[
 q\le G\le q^2-q,
 \qquad0<(M_C+q\MF)J<2q^2-q-1.
\]
The native mask bounds used here are \(0<M_C,M_F<B-1\), with
\(\MF=M_F+B-1\). Consequently
\begin{equation}\label{eq:Rbounds}
 q(q^2-1)<R<q^4-q^3+q^2-1<q^4-2,
 \qquad E\ge q^4>R+2,\qquad a>R+2.
\end{equation}
These bounds precede every dyadic, Pell-index, or computation-decoding
conclusion.

\subsection{Bootstrap, exact index, and the two remaining signs}

Put \(A=a+2\) and \(P_0=2XY^2+1\). The first norm gives
\[
 |\tau|=\chi_{P_0}(n),\qquad k=2\psi_{P_0}(n),\qquad n\ge1.
\]
Since \(P_0\equiv1\pmod E\), the index unit gives
\[
 2n\equiv R+\varepsilon\pmod E,
 \quad 2n=R+\varepsilon+vE\quad(v\ge0),
 \quad n\ge(R-1)/2\ge24.
\]
This does not require a sign for \(h\). The main norm gives
\(|D|=\chi_A(p)\), \(c=\psi_A(p)\). Since \(P_0>A\) and \(c>k\),
monotonicity gives \(p>n\). Elementary growth now implies
\[
 c>A\Delta^2,\qquad c>2p,\qquad c>2R.
\]
For the last inequality one may already use \(p\ge2\) and
\(c\ge2A>2R\); no positivity of \(h\) is hidden in it.

Lemma~\ref{lem:rank} applies to \(I=ic^2\), yielding the auxiliary
index \(m\), its divisibility, and \(c^2\mid f^2-1\).
The formula for \(V\), even when \(T\) and \(f\) are signed, gives
\[
 V\equiv-c\pmod{|f|},\qquad V\equiv-R\pmod c.
\]
Lemma~\ref{lem:signed} therefore gives \(p=R\).

If \(v\ge1\), then \(2n\ge R-1+E>2R\), contrary to \(n<p=R\).
Thus \(2n=R+\varepsilon\). If \(\varepsilon=-1\), then
\(p=2n+1\). Let \(Q=\chi_A(2)=2A^2-1>P_0\). Duplication gives
\[
 c=\psi_A(2n+1)>\psi_A(2n)
  =2A\psi_Q(n)\ge2A\psi_{P_0}(n)=Ak>k(Y+1),
\]
contrary to \(kY<c<k(Y+1)\). Hence \(\varepsilon=\nu=1\).
In the product presentation this now excludes \(C=0\), so \(C>0\).
We have \(R=2n-1\), an odd positive integer.

\subsection{The main root sign and the power of two}

The sign of \(D\) matters if \(\gamma\) is unrestricted. Write
\(R=2r+1\), \(r=n-1\). The actual lower bound in
\eqref{eq:Rbounds} gives \(r\ge24\). Standard recurrence bounds give
\[
 \psi_b(t)\ge(2b-1)^{t-1},\qquad
 \psi_b(t)\le(2b)^{t-1}\quad(b\ge2,t\ge1).
\]
Using \((2A-1)^2>2P_0X\), the strict ratio implies
\[
 Y+1>\frac{c}{k}
 =\frac{\psi_A(2r+1)}{2\psi_{P_0}(r+1)}>\frac{X^r}{2}.
\]
Thus
\begin{equation}\label{eq:Ylarge}
 Y>X^r/3>2^R,\qquad Y>X,\qquad a>2^R.
\end{equation}
The inequalities use only \(X\ge16\) and \(r\ge24\).

Modulo the odd integer \(H=4a+3\), the two projected Pell sequences
obey
\[
 \chi_A(t)-a\psi_A(t)\equiv2^t,
 \qquad\chi_A(t)+a\psi_A(t)\equiv2^{-t}.
\]
One checks the initial values and recurrence; \(2(a+1)\) is the
inverse of 2 modulo \(H\). If \(D<0\), its defining equation
would give \(X2^R\equiv-1\pmod H\). But
\(0<X2^R<XY=E<a<H-1\), a contradiction. Therefore \(D>0\).
The first projected congruence now gives \(X\equiv2^R\pmod H\).
Both are between 0 and \(a<H\), hence
\begin{equation}\boxed{X=2^R.}\label{eq:power}\end{equation}
In particular \(q\) is a power of two, and \((B-1)\mid(q-1)\)
implies \(q=B^N\) for an integer \(N\ge1\).

Also \(h>0\) and \(\gamma>0\) are consequences, not domain
assumptions. A negative \(h\) contradicts \(k>0\) and
\(E>R+2\); \(h=0\) contradicts first-norm growth. Finally
\[
 D-ac=2c-\psi_A(R-1)>c>Y>X,
\]
so the equation defining \(D\) forces \(\gamma>0\).

\subsection{The inherited bit-count and native compiler interfaces}

At this point the first/main data satisfy the hypotheses of the
post-index ratio and valuation part of the half-binomial Pell
argument~\cite{halfbinomial,asymmetric}: its scale assumptions are
\(X=wq\), \(Y=sq^3\), \(q\ge16\), and
\(3q+1\le R<q^4\), with the exact odd index and strict ratio just
proved. Its input-independent conclusion is
\[
 \vv(Y)=\pc(R)-2.
\]
Here is the precise tail of the argument being used. Put
\(\xi=(X+1)^{2r}/X^r\). The inherited ratio estimate gives
\[
 0<\frac{c}{k}-\frac{\xi}{2}
   <\frac{16r}{X+1}<\frac12.
\]
Expand \(\xi=M+\theta\), where
\[
 M=\binom{2r}{r}+
       \sum_{j=1}^{r}\binom{2r}{r+j}X^j,
 \qquad 0<\theta<\frac14.
\]
The integer \(M\) is even. Thus \(Y<c/k<Y+1\) gives
\(Y=M/2\). The central-binomial valuation identity
\(\vv\binom{2r}{r}=\pc(r)\), together with \(X=2^{2r+1}\),
now gives \(\vv(Y)=\pc(r)-1=\pc(R)-2\).
This uses no auxiliary positivity or positive-root parity assertion.
If the symmetric version is used, \(R>3q\) and \(X=2^R\) imply
\(q^3\mid X\), restoring its scale by a positive integer rescaling
of \(w\). Since \(q^3\mid Y\), writing \(q=2^t\) gives
\(\pc(R)\ge3t+2\).

The shifted packing identity is
\begin{align*}
 S'&=C+qF-1,\\
 T'&=(M_CJ+1)+q(M_FJ-1),\\
 R&=(q^2-S')(q^2-1)+T'.
\end{align*}
Here \(0<S'<q^2\) and the two mask words have the precise native
populations. The inverse-population lemma~\cite{raw75,blank06} gives
\[
 (C-1)\mathbin{\&}(M_CJ+1)=0,
 \qquad F\mathbin{\&}(M_FJ-1)=0.
\]
These are decoded consequences, not operations in the polynomial.
Transport is the required native congruence because
\((K+X)C-F\) differs from \((K+w)C-F\) by a multiple of \(q-1\).
Sections~\ref{sec:local}--\ref{sec:coefficients} establish the
remaining finite-language and coefficient hypotheses. The masks force
a Start-containing valid blank tableau, hence a halting K215 run.

The half-binomial, inverse-population, and canonical actual-index
completeness lemmas are inherited mathematical dependencies. Their
statements, pinned files, and relevant source proofs are included or
identified in \code{provenance.json}. The finite scripts in this
archive are not substitutes for those unbounded arguments.

\subsection{Completeness and the positive transport shift}

Given a blank halting run, the native compiler's canonical padding and
actual-index construction supplies the normalized blank witnesses of
package 06~\cite{blank06}. Keep every coordinate except the ordinary
strong multiplier and set \(i_{\rm ordinary}=\Delta i_{\rm normalized}\).
Then
\[
 (i_{\rm ordinary}c^2)^2
 =\Delta(f^2-1),
\]
and the auxiliary coefficient has the same value. Thus the ordinary
system with the old transport factor has a positive solution.

On this solution, \(w=2^R/q>q>F\), because \(R>3q\).
As \(C\ge1\),
\[
 (z_0-1)(q-1)=(K+w)C-F>0.
\]
Hence \(z_0\ge2\), and \(z=z_0-1\) is positive. This is the
one-operation transport saving in \eqref{eq:system60}. The reverse
shift is immediate. Completeness on the larger signed domain follows
because the canonical witnesses are positive.

\section{Fourteen unknowns by eliminating the strong root}
\label{sec:fourteen}
The blank-input reduction also combines with the root-elimination method
of packages 03 and 05. This removes one supplied unknown, at the cost
of a higher-degree equation. In this section \(i\) has its
\emph{normalized} meaning: the root to be eliminated would satisfy
\(f^2=1+\Delta i^2c^4\). The ordinary multiplier used in
Section~\ref{sec:arithmetic-proof} will be \(\Delta i\).

Delete \(f\) from the fifteen-variable list and keep all definitions
that do not involve it. Use the old positive transport coordinate
\(z_0\), so that \(N_t=(K+w)C+q-F-z_0(q-1)\). Define
\begin{align}
 Q_*&=(\Delta ic^2)^2,\qquad d_*=1+\Delta i^2c^4,\notag\\
 b_*&=c+Rd_*,\qquad v_*=cT,\notag\\
 A_*&=Q_*(d_*v_*^2+b_*^2-y^2)+y^2-1,\notag\\
 B_*&=2Q_*v_*b_*,\qquad K_*=A_*^2-d_*B_*^2,\notag\\
 P_4&=N_0N_1N_kN_t.\label{eq:rootdata14}
\end{align}
All these are polynomial abbreviations in fourteen supplied unknowns;
none introduces another existential coordinate. The first equation is
\begin{equation}\boxed{\Phi_{14}=P_4(1+K_*)-1=0.}
\label{eq:phi14}\end{equation}
Its circuit has 75 operations and exact degree 293.

A second finalizer incorporates the outer product inside the norm.
Put
\begin{align*}
 A_{\mathrm{out}}
   &=P_4\bigl[Q_*(d_*v_*^2+b_*^2)+(1-Q_*)y^2\bigr]-1,\\
 B_{\mathrm{out}}&=P_4B_*.
\end{align*}
Then
\begin{equation}\boxed{\Psi_{14}
 =A_{\mathrm{out}}^2-d_*B_{\mathrm{out}}^2=0}
\label{eq:psi14}\end{equation}
has 73 operations and exact degree 342. Both equations initially use
fourteen strictly positive unknowns. Their coefficients are the same
four fixed integers as before.

The circuits compute the norm cross-term through
\[
 d_*B_*^2=4(Q_*d_*v_*^2)(Q_*b_*^2),
\]
and its analogous outer-product expression. Thus they do not pay to
construct \(B_*\) separately. Neither circuit contains division,
square-root extraction, variable exponentiation, or a free program index.

\subsection{Restoring an integer root, with its sign left free}

At a zero of \eqref{eq:phi14}, both \(P_4\) and \(1+K_*\) are
integer units. Because \(B_*\) is even,
\[
 1+K_*\equiv1+A_*^2\pmod4,
\]
which cannot be \(-1\). Hence \(K_*=0\) and \(P_4=1\).
Every factor inside \(P_4\) is consequently a unit. The first and
main norms equal \(1\) by the descent and congruence arguments already
given. The transport-unit argument yields \(C\ge0\), and the
pre-decoding bounds \eqref{eq:Rbounds} apply; in particular \(R>0\).
On the positive domain, \(Q_*,v_*,b_*>0\), so \(B_*\ne0\).

The rational number
\[
 f=\frac{A_*}{B_*}
\]
satisfies \(f^2=d_*\). A rational square root of an integer is an
integer: write it in lowest terms, and any prime in its denominator
would have negative valuation in its square. Since \(d_*>1\),
this restored integer \(f\) is nonzero. Its sign need not be determined.
Set \(V=v_*f-b_*\). Using \(f^2=d_*\), the exact remainder identity is
\begin{equation}
 A_*-B_*f=Q_*(V^2-y^2)+y^2-1=0.
 \label{eq:root-remainder14}
\end{equation}
Moreover
\[
 V=c(Tf-1)-Rf^2,\qquad f^2-\Delta i^2c^4=1.
\]
Thus an integer solution of the normalized blank equations has been
restored, allowing a signed \(f\). Replace the multiplier by
\(i_{\mathrm{ordinary}}=\Delta i>0\). Then
\[
 (i_{\mathrm{ordinary}}c^2)^2
   =\Delta(f^2-1)=Q_*.
\]
All factors of \eqref{eq:p69} have the required product, so the
signed-domain theorem of Section~\ref{sec:arithmetic-proof} recovers
a genuine K215 blank halting run. No separate positivity theorem for
the reconstructed root is required.

For \eqref{eq:psi14}, first note that \(P_4=0\) would make its
left side equal \(1\). Otherwise reduction modulo \(P_4\) gives
\(\Psi_{14}\equiv1\pmod{P_4}\), so a zero forces
\(P_4=\pm1\). Its four factors are units and give the same
outer bounds. Since \(B_{\mathrm{out}}\ne0\), set
\(f=A_{\mathrm{out}}/B_{\mathrm{out}}\). Again \(f^2=d_*\)
makes \(f\) an integer. The corresponding remainder identity is
\[
 A_{\mathrm{out}}-B_{\mathrm{out}}f
  =P_4\bigl[Q_*(V^2-y^2)+y^2\bigr]-1=0.
\]
After the same change of strong multiplier, this is exactly the
product needed by \eqref{eq:p69}. The argument does not assume the
sign of \(P_4\) in advance.

Conversely, a canonical positive blank compiler solution has
\(f^2=d_*\), \(A_*=B_*f\), and \(P_4=1\). Deleting \(f\)
therefore makes both \eqref{eq:phi14} and \eqref{eq:psi14} vanish.
This proves the two directions of their existence equivalence and
hence their stated \(\Con(\mathrm{ZFC})\) equivalences.

\subsection{A signed extension and 32 unrestricted integer unknowns}

The nine coordinates
\[
 J,F,\alpha,z_0,i,s,w,\eta,\zeta
\]
can remain positive while the other five,
\(h,T,\tau,y,\gamma\), are arbitrary integers. The only new
obstruction to the preceding argument would be \(T=0\), making
the reconstructed-root denominator vanish. We exclude it before
taking a quotient.

For either polynomial the outer unit argument already gives
\(0<R<a<\Delta\) and \(c>2\). Since
\(Q_*=\Delta(d_*-1)\),
\begin{align*}
 Q_*-b_*&=(\Delta-R)d_*-\Delta-c\\
 &\ge d_*-\Delta-c
   =1+\Delta(i^2c^4-1)-c>0.
\end{align*}
Consequently \(2<b_*<Q_*\). If \(T=0\), then \(v_*=B_*=0\).
For \(\Phi_{14}=0\), the equation \(K_*=0\) forces \(A_*=0\),
and therefore
\begin{equation}
 Q_*b_*^2-(Q_*-1)y^2=1.\label{eq:zero-T-norm}
\end{equation}
For \(\Psi_{14}=0\), the equation
\(A_{\mathrm{out}}=0\) says that the left side of
\eqref{eq:zero-T-norm} equals \(P_4=\pm1\). It cannot equal
\(-1\), since \(Q_*\) is a square and this expression is a
square modulo four. Thus \eqref{eq:zero-T-norm} follows again.

Every solution of this equation with \(b_*>1\) has
\(b_*\ge2Q_*-1\). Indeed,
\[
 b_*^2-1=(Q_*-1)(y^2-b_*^2)
\]
implies \(|y|\ge b_*+1\), and for \(1<b_*\le2Q_*-2\)
one has \(b_*^2-1<(Q_*-1)(2b_*+1)\), a contradiction.
This conflicts with \(b_*<Q_*\). Hence \(T\ne0\), the
root denominators are nonzero, and the signed reconstruction is valid.
The signed ordinary theorem then supplies soundness exactly as above.

Use the positivity polynomial \eqref{eq:pos} independently for the
nine guarded coordinates and leave the other five unrestricted.
The resulting polynomials have \(9\cdot3+5=32\) integer unknowns.
The conversion is surjective onto the required positive values, so it
preserves existence in both directions. Their operation counts and
degrees are included below. They improve the variable count of the
explicit family studied here; they do not assert a global minimum over
all Diophantine representations.

\subsection{Exact degrees and operation ledgers}

Write \(Q_0=\beta J\), \(k_0=\eta+\zeta\), and
\(T_2=w(Q_0-F-\alpha)-z_0Q_0\). The two leading homogeneous
forms are respectively
\begin{align}
 (\Phi_{14})_{293}
  &=8h\gamma i^{12}k_0^{27}w^{21}s^{48}
       Q_0^{177}(Q_0-F)^4T_2,\label{eq:phi14-leader}\\
 (\Psi_{14})_{342}
  &=64h^2\gamma^2 i^{12}k_0^{30}w^{26}s^{56}
       Q_0^{206}(Q_0-F)^4T_2^2.\label{eq:psi14-leader}
\end{align}
They are nonzero for every \(\beta>0\). For example, the coefficients
at
\[
 J^{182}h\gamma i^{12}\eta^{27}w^{21}s^{48}z_0,
 \qquad
 J^{212}h^2\gamma^2i^{12}\eta^{30}w^{26}s^{56}z_0^2
\]
are \(-8\beta^{182}\) and \(64\beta^{212}\), respectively.

To check the degree accounting, \(Q_*,d_*,b_*,v_*\) have degrees
\(46,34,38,6\). The unique top term of \(A_*\) is \(Q_*b_*^2\),
of degree 122. Hence \(K_*\) has degree 244: the competing
term \(d_*B_*^2\) has degree 214. The degree-49 outer product gives
\(244+49=293\). In the second construction the squared leading
term has degree 342 and its subtracted cross-term degree 312.

For the integer versions assign weight two to the nine guarded
coordinates and weight one to the five signed coordinates. The leading
forms above have signed weights two and four. Their exact degrees
therefore become \(2\cdot293-2=584\) and
\(2\cdot342-4=680\). Direct weighted upper bounds give the same
values. Disjoint quadratic positivity substitutions preserve
nonvanishing, as explained in Section~\ref{sec:integer}.

\begin{center}\small
\begin{tabular}{@{}lrrrr@{}}
\toprule
Presentation & Domain & Unknowns & Degree & Operations\\
\midrule
\(\Phi_{14}\) & positive &14&293&75 = 44M+31A\\
\(\Psi_{14}\) & positive &14&342&73 = 44M+29A\\
\(\Phi_{14}\) after positivity substitution & integers &32&584&138\\
\(\Psi_{14}\) after positivity substitution & integers &32&680&136\\
\bottomrule
\end{tabular}
\end{center}
The integer conversion adds seven gates for each of nine positivity
substitutions. The files \code{arithmetic/build_fourteen.py} and its
emitted packets give complete evaluators and reproduce the counts,
root-remainder identities, substitution identities, and degree
diagnostics. These finite checks supplement the proofs; they do not
replace the unbounded compiler or Pell arguments.


\section{Lower degree through an auxiliary gap}
\label{sec:gap15}
A further change of coordinate lowers the degrees of the ordinary
construction. Replace the supplied \(y\) by a supplied \(e\), and
compute \(y=V+e\). The auxiliary factor becomes
\begin{equation}
 N_A^{\mathrm{gap}}
  =U\bigl(V^2-(V+e)^2\bigr)+(V+e)^2
  =V^2-(U-1)e(2V+e).
 \label{eq:gap-auxiliary}
\end{equation}
Keep the other factors and residuals unchanged. Thus the new product is
\[
 P_{70}^{\mathrm{gap}}
   =N_0N_1N_A^{\mathrm{gap}}N_kN_tN_s-1=0.
\]
and the new six-equation system is
\begin{equation}\boxed{\begin{aligned}
 N_0=N_1=N_A^{\mathrm{gap}}=N_k&=1,\\
 (K+w)C&=F+z(q-1),\\
 I^2&=U.
\end{aligned}}\label{eq:gap-system61}\end{equation}
It replaces only the third equation of
\eqref{eq:system60} by \(N_A^{\mathrm{gap}}=1\).
The sum of its six squared residuals gives a third presentation.
The circuit needs just one additional addition: its existing
\(y\)-square now squares the computed value \(V+e\).

The signed-domain equivalence is immediate from the inverse integer
substitutions \(e=y-V\) and \(y=V+e\). The nine coordinates
\(J,F,\alpha,z,i,s,w,\eta,\zeta\) remain positive, while
\(h,f,T,\tau,e,\gamma\) may be arbitrary integers. Soundness therefore
follows from Section~\ref{sec:arithmetic-proof}. Positive-domain
completeness needs one extra check. The canonical positive auxiliary
Pell point used by the blank compiler has \(V,y>0\), \(U>1\),
\(c>2R\), and \(V\equiv-R\pmod c\). Hence
\(V\ge c-R>1\), and its norm equation gives
\[
 (U-1)(y^2-V^2)=V^2-1>0.
\]
Thus \(y>V\) and the canonical gap \(e=y-V\) is strictly positive.
This proves the same existence equivalence in both stated domains.

\begin{center}
\begin{tabular}{lrrr}
\toprule
Presentation & Unknowns & Operations & Exact degree\\
\midrule
Six equations & \(15\) positive & 61 & maximum 22\\
Product & \(15\) positive & 70 & 93\\
Sum of squares & \(15\) positive & 78 & 44\\
Six equations & \(33\) integer & 124 & maximum 44\\
Product & \(33\) integer & 133 & 179\\
Sum of squares & \(33\) integer & 141 & 88\\
\bottomrule
\end{tabular}
\end{center}

For exact degree certificates, set \(Q_0=\beta J\),
\(k_0=\eta+\zeta\), and
\(T_2=w(Q_0-F-\alpha)-z_0Q_0\). The auxiliary leader is
\(-2ek_0w^2s^3Q_0^{11}Tf^3\); the six factor degrees are
\(22,18,22,7,2,22\). Consequently
\begin{equation}
 (P_{70}^{\mathrm{gap}})_{93}
  =-16Q_0^{52}h\gamma i^2k_0^8w^7s^{15}Tf^3eT_2.
 \label{eq:gap-product-leader}
\end{equation}
The coefficient of
\(J^{53}h\gamma i^2\eta^8w^7s^{15}Tf^3ez_0\)
is \(16\beta^{53}\ne0\). The degree-22 leaders of the first,
auxiliary and strong residuals are
\[
 -w^2k_0^2s^4Q_0^{14},\qquad
 -2ek_0w^2s^3Q_0^{11}Tf^3,\qquad
 i^2k_0^4s^4Q_0^{12}.
\]
Their squared sum is nonzero, proving degree 44 for the sum of squares.

For the integer conversion, give weight two to the nine positive
coordinates and weight one to the other six. The residual weights are
\(44,35,39,13,4,44\). Nine independent substitutions
\eqref{eq:pos} add 63 gates and give 33 integer unknowns. The product
has degree 179, while the first and strong squared residuals give
degree 88 for the sum of squares. The disjoint homogeneous quadratic
substitutions preserve the nonzero leading forms. Complete circuits,
ledgers and checks are emitted by \code{arithmetic/build_gap.py}.


\section{The degree-219 equation in fourteen unknowns}
\label{sec:gap14}
The auxiliary gap can be introduced \emph{before} eliminating the
normalized strong root. This improves the degree-293 inner norm in
Section~\ref{sec:fourteen} to degree 219 at the same 75-operation
cost. It still uses fourteen positive unknowns. Keep normalized \(i\),
delete \(f\), replace \(y\) with \(e\), and define
\begin{align}
 Q_*&=(\Delta ic^2)^2,& d_*&=1+\Delta i^2c^4,\notag\\
 b_*&=c+Rd_*,&v_*&=cT,\notag\\
 S_*&=Q_*-1,&r_*&=b_*+S_*e,\notag\\
 p_*&=d_*v_*^2,&B_*^{\mathrm{gap}}&=2v_*r_*,\notag\\
 A_*^{\mathrm{gap}}
   &=p_*+r_*^2-Q_*S_*e^2-1,\notag\\
 K_*^{\mathrm{gap}}
   &=(A_*^{\mathrm{gap}})^2-d_*(B_*^{\mathrm{gap}})^2,
   &P_4&=N_0N_1N_kN_t.
 \label{eq:gap-root-data}
\end{align}
Here the four factors use the old transport coordinate \(z_0\).
The proposed equation is
\begin{equation}\boxed{
 \Phi_{14}^{\mathrm{gap}}
    =P_4(1+K_*^{\mathrm{gap}})-1=0.}
 \label{eq:gap14}
\end{equation}
Its evaluator computes the cross-term as \(4p_*r_*^2\), so
\(B_*^{\mathrm{gap}}\) is a proof abbreviation and needs no separate
register. Removing \(f\) and the old \(y\)-square leaves 49 prefix
operations; the new tail costs 26. The ledger is 43 multiplications
and 32 additions or subtractions.

The degree improvement comes from the exact cancellation
\begin{equation}
 A_*^{\mathrm{gap}}
  =p_*+b_*^2+(Q_*-1)e(2b_*-e)-1.
 \label{eq:gap-root-cancellation}
\end{equation}
If \(f^2=d_*\) and \(V=v_*f-b_*\), then
\begin{equation}
 A_*^{\mathrm{gap}}-B_*^{\mathrm{gap}}f
  =Q_*\bigl[V^2-(V+e)^2\bigr]+(V+e)^2-1.
 \label{eq:gap-root-remainder}
\end{equation}
Thus the new norm eliminates precisely the gap auxiliary equation.

\subsection{Restoration and the signed-domain denominators}

At a zero of \eqref{eq:gap14}, both \(P_4\) and
\(1+K_*^{\mathrm{gap}}\) are integer units. Because
\(B_*^{\mathrm{gap}}\) is even, the latter is congruent to
\(1+(A_*^{\mathrm{gap}})^2\) modulo four and cannot be \(-1\).
Hence \(P_4=1\) and \(K_*^{\mathrm{gap}}=0\). The outer unit
arguments already give \(N_0=N_1=1\) and
\(0<R<a<\Delta\). On the positive domain,
\(v_*,b_*,r_*>0\), so the denominator in
\(f=A_*^{\mathrm{gap}}/B_*^{\mathrm{gap}}\) is nonzero.
This rational root satisfies \(f^2=d_*\), hence is an integer.
Identity~\eqref{eq:gap-root-remainder}, with \(y=V+e\), restores
the normalized auxiliary equation. Replacing the multiplier by
\(i_{\mathrm{ordinary}}=\Delta i\) gives

\[
 (i_{\mathrm{ordinary}}c^2)^2
   =\Delta(f^2-1)=Q_*.
\]
The signed ordinary theorem therefore proves soundness without imposing
a sign on the recovered \(f\) or \(y\).

For the expanded domain, keep only
\(J,F,\alpha,z_0,i,s,w,\eta,\zeta>0\), and allow
\(h,T,\tau,e,\gamma\in\Z\). We must now exclude two possible
zero factors in \(B_*^{\mathrm{gap}}=2v_*r_*\).
The stronger version of the earlier outer bound is
\begin{align*}
 Q_*-b_*&=(\Delta-R)d_*-\Delta-c\\
 &\ge1+\Delta(i^2c^4-1)-c>1.
\end{align*}
Indeed \(c>2\), \(i\ge1\), and \(\Delta>0\) are already known.
Thus \(2<b_*<Q_*-1=S_*\), so

\[
 r_*=b_*+S_*e\ne0
 \qquad\text{for every integer }e.
\]

If \(T=0\), then \(v_*=p_*=B_*^{\mathrm{gap}}=0\), and
\(K_*^{\mathrm{gap}}=0\) forces \(A_*^{\mathrm{gap}}=0\).
Equivalently,
\begin{equation}
 Q_*b_*^2-(Q_*-1)(e-b_*)^2=1.
 \label{eq:gap-zero-T}
\end{equation}
Writing \(y_0=e-b_*\), the inequality \(b_*>1\) would imply
\(|y_0|\ge b_*+1\), and hence
\[
 b_*^2-1=(Q_*-1)(y_0^2-b_*^2)
       \ge(Q_*-1)(2b_*+1).
\]
This contradicts \(1<b_*<Q_*\). Therefore \(T\ne0\), both
denominator factors are nonzero, and the signed reconstruction is valid.

For completeness, the canonical positive blank witness satisfies
\(f^2=d_*\), \(P_4=1\), and \(e=y-V>0\), as proved for the
gap auxiliary coordinate. Equation~\eqref{eq:gap-root-remainder}
then gives \(A_*^{\mathrm{gap}}=B_*^{\mathrm{gap}}f\), so deleting
\(f\) yields a positive zero of \eqref{eq:gap14}. This proves
both directions of the same blank-halting and
\(\Con(\mathrm{ZFC})\) equivalence.

Nine independent positivity substitutions \eqref{eq:pos}, with the
other five coordinates left unrestricted, give a polynomial in
\(9\cdot3+5=32\) integer unknowns and 138 arithmetic operations.

\subsection{Exact degrees and the outer-norm comparison}

For the degree calculation, the following table also records weights
two on the nine positive coordinates and one on the five signed ones.
\begin{center}
\begin{tabular}{lrrrrrrrrr}
\toprule
Quantity & \(Q_*\)&\(d_*\)&\(b_*\)&\(v_*\)&\(p_*\)&\(r_*\)
 &\(A_*^{\mathrm{gap}}\)&\(B_*^{\mathrm{gap}}\)&\(P_4\)\\
\midrule
Degree &46&34&38&6&46&47&85&53&49\\
Weight &92&68&76&11&90&93&169&104&96\\
\bottomrule
\end{tabular}
\end{center}
The unique highest term of \eqref{eq:gap-root-cancellation} is
\(2Q_*eb_*\). Its square has degree 170 and weight 338, while
\(d_*(B_*^{\mathrm{gap}})^2\) has degree 140 and weight 276.
Consequently the new equation has exact degree \(49+170=219\),
and its integer conversion has exact degree \(96+338=434\).

Using \(Q_0=\beta J\), \(k_0=\eta+\zeta\), and
\(T_2=w(Q_0-F-\alpha)-z_0Q_0\), the leading form is
\begin{equation}
 (\Phi_{14}^{\mathrm{gap}})_{219}
   =32h\gamma e^2i^8k_0^{19}w^{17}s^{36}
       Q_0^{131}(Q_0-F)^2T_2.
 \label{eq:gap-root-leader}
\end{equation}
Its coefficient at
\(J^{134}h\gamma e^2i^8\eta^{19}w^{17}s^{36}z_0\)
is \(-32\beta^{134}\ne0\). This is also the weighted leading form
before the homogeneous positivity substitution, which preserves
nonvanishing.

For comparison, the outer-norm variant is
\[
 A_{\mathrm{out}}^{\mathrm{gap}}
   =P_4(A_*^{\mathrm{gap}}+1)-1,
 \qquad
 B_{\mathrm{out}}^{\mathrm{gap}}=P_4B_*^{\mathrm{gap}},
\]
\[
 \Psi_{14}^{\mathrm{gap}}
   =(A_{\mathrm{out}}^{\mathrm{gap}})^2
      -d_*(B_{\mathrm{out}}^{\mathrm{gap}})^2=0.
\]
It also costs 75 operations, but has degree 268; its 32-integer-variable
conversion has 138 operations and degree 530. The leading form is
\[
 256h^2\gamma^2e^2i^8k_0^{22}w^{22}s^{44}
       Q_0^{160}(Q_0-F)^2T_2^2.
\]
This valid form is dominated by \eqref{eq:gap14} at the same gate
count. It does not dominate the earlier 73-operation outer norm.

Its soundness proof is included to make that comparison independently
checkable. Reduction modulo \(P_4\) forces \(P_4=\pm1\), giving
the same outer bounds and \(r_*\ne0\). If \(T=0\), its equation says
that \(P_4\) times the left side of \eqref{eq:gap-zero-T} equals one.
That bracket cannot be \(-1\) modulo four because \(Q_*\) is a square;
therefore \eqref{eq:gap-zero-T} holds, an impossibility. Its root
denominator is consequently nonzero. The integer restored root now
satisfies
\[
 A_{\mathrm{out}}^{\mathrm{gap}}
    -B_{\mathrm{out}}^{\mathrm{gap}}f
  =P_4\bigl\{Q_*[V^2-(V+e)^2]+(V+e)^2\bigr\}-1=0.
\]
This restores the full normalized product, hence the signed ordinary
one. The canonical witness proves completeness as before.

Both gap root variants, their integer conversions, their live-gate
ledgers, and their algebraic checks are emitted by
\code{arithmetic/build_gap.py}. The separate proof note records the
complete coordinate substitutions and nonzero-denominator arguments.


\section{A quadratic system in thirty-two unknowns}
\label{sec:quartic}

Retain the same fixed coefficients. Use the 32 positive unknowns
\begin{gather*}
 J,F,\alpha,z,f,h,i,s,w,\tau,\eta,\zeta,y,\gamma,o,j,\\
 N,E,P,c,a,X,D,C_2,G,L,F_2,V,g,I,B_a,\Delta.
\end{gather*}
Abbreviate \(q=\beta J+1\), \(k=\eta+\zeta\), \(H=4a+3\),
\(C=q-F-\alpha\), and \(R=k-1-hE\). These are substitutions, not
new unknowns. The following 23 equations are quadratic after expansion:
\begin{align}
 L&=q^2-1,& N&=q(L+1),&X&=wq,\notag\\
 a&=E+sN,&E&=X(a-E),&c&=k(a-E)+\eta,\notag\\
 P&=E(c-\eta)+k,&\tau^2&=P(P-k)+1,\notag\\
 D&=X+ac+\gamma H,&\Delta&=a^2+H,&C_2&=c^2,\notag\\
 D^2&=\Delta C_2+1,&G&=(q-1)(q-F)+\alpha,\notag\\
 R&=LG+(M_C+q\MF)J,&(K+w)C&=F+z(q-1),\notag\\
 F_2&=f^2,&V&=of-c,&V&=jc-R,\notag\\
 g&=y^2-V^2,&I&=iC_2,&B_a&=I^2,\notag\\
 B_a&=\Delta(F_2-1),&y^2&=B_ag+1.
 \label{eq:quadratic32}
\end{align}
The expanded residuals are supplied in \code{quadratic/quartic32.json}.
Taking the sum of their squares gives an exact quartic: only the first
norm residual contains \(\tau\), and its square contributes
\(\tau^4\) with coefficient 1.

\begin{proposition}
The positive solvability of \eqref{eq:quadratic32} is equivalent to
that of \eqref{eq:system60}.
\end{proposition}
\begin{proof}
The first equations give \(Y=a-E=sq^3\), \(E=XY\),
\(c=kY+\eta\), and \(P=EkY+k\). Thus its first norm is exactly
\(\tau^2-XY^2(XY^2+1)k^2=1\). The main, ordinary strong, packing,
and transport equations are the same. The last equations give
\(I^2V^2-(I^2-1)y^2=1\) with \(V=of-c=jc-R\).

The ordinary rank lemma gives \(c^2\mid f^2-1\). Set
\(T=(o+Rf)/c\). This is integral because
\[
 o+Rf=f(of+R)-o(f^2-1),\qquad of+R=(j+1)c.
\]
It is positive, and \(c(Tf-1)-Rf^2=of-c=V\), recovering the
fifteen-coordinate system.

Conversely, use a canonical positive solution, or apply the positive
auxiliary restoration argument of package 02. The expressions
\[
 o=cT-Rf,\qquad j=Tf-1-R(f^2-1)/c
\]
are positive integers. For clarity, positivity follows from
\(U>Rf^2+c\), \(V>-U\), the auxiliary norm's gap
\(|V|>1\Rightarrow |V|\ge2U-1\), and the exclusion of
\(V=0,\pm1\) by \(f>2c\) and \(V\equiv-c\pmod f\).
Thus \(V>1\). Since \((U-1)(y^2-V^2)=V^2-1\), the supplied
\(g=y^2-V^2\) is positive. All other introduced registers have
positive values by their displayed definitions.
\end{proof}

The ordinary strong norm removes the old intermediate \(\Delta c^2\)
from package 08's lift. Restoring \(o,j\) instead of \(T,Tf-1\)
then makes \(R\) occur only linearly, allowing its elimination by
\(R=k-1-hE\). These two changes account for the improvement from
34 variables/25 quadratics to 32 variables/23 quadratics.

\section{The local language and blank-halting equivalence}
\label{sec:local}

The fixed machine is \code{candidate-k215/k215.tm}, whose SHA-256 is
\begin{center}\footnotesize
\code{fdbcef0825886035c3091c8537a69535bc96f9e265f0ae45863543d40ae70090}.
\end{center}
It has 215 working states, two tape symbols, a blank two-sided tape, start
state 0, and a separate halt 215. The first transition is
\((0,0)\mapsto(1,+1,93)\); the unique entering-halt transition is
\((68,1)\mapsto(1,+1,215)\). For tableau padding only, append
\((215,b)\mapsto(b,0,215)\).

The 439 cell symbols are ordinary symbols 0 and 1; headed symbols
\(2(q_0+1)+b\) for \(0\le q_0\le215\), \(b\in\{0,1\}\);
initialization symbols \(L,Q,R\) with codes 434,435,436; vertical
boundary 437; and horizontal boundary 438. Initialization payloads are
blank, with the start head at \(Q\).

The package-06 center predicate uses only the prior, left, center, right,
and following cells. It propagates the vertical-boundary bit in time;
forbids adjacent vertical boundaries and consecutive horizontal boundary
rows; makes horizontal boundaries uniform inside a strip; and makes the
initial row exactly \(L^+QR^+\). Away from boundaries it applies the
literal three-neighbor machine update, forbidding multiple head arrivals
and escaping heads. The last interior row may contain a head only in
the halt state. The complete predicate is preserved unchanged in
\code{compiler/blank_helical.py}; its boundary clauses and exact
machine table are part of the finite coefficient definition.

\subsection{Why five positions suffice}

A cross is \((u,l,c,r,d)\), in prior/left/center/right/following order.
It is locally valid if the old nine-cell predicate accepts it with
ordinary blank corners. This is equivalent to every choice of corners,
because the predicate never reads them. The old selector count is
therefore exactly \(439^4\) times the cross count.

Enumerate the following two crosses first:
\[
 S=(438,434,435,436,1),\qquad
 E_*=(438,0,437,438,437).
\]
Then enumerate all other valid crosses lexicographically. \(S\) is
the initialized start cross; its following center is the ordinary
symbol 1 written by the first transition. \(E_*\) is locally valid
at its vertical-boundary center, but cannot occur globally: the prior
horizontal-boundary cell has vertical bit 0, whereas its following
center has vertical bit 1. Its exclusion removes no genuine tableau.

The exact number of valid crosses is
\begin{equation}\label{eq:cross-count}
 \kappa=38\,256\,607\,897.
\end{equation}
The inherited thirteen-category census counts these crosses directly;
the new audit replays each category against the exact predicate.

Let the right neighbor of cyclic cell \(i\) be \(i-1\), and its
following-time neighbor be \(i-h\). Four equality families suffice:
\[
 r_i=c_{i-1},\qquad c_i=l_{i-1},\qquad
 d_i=c_{i-h},\qquad c_i=u_{i-h}.
\]
Applying them at neighboring indices gives
\begin{equation}\label{eq:reconstruct}
 W_i=(c_{i+h},c_{i+1},c_i,c_{i-1},c_{i-h}).
\end{equation}
This holds even when cyclic indices coincide. Pullback along
\((x,t)\mapsto-x-ht\pmod N\) therefore reconstructs a periodic
plane satisfying the exact center predicate.

\subsection{A marked periodic plane contains a finite halting run}

At a Start, scan its initialization row to the left and right. The
\(L\)-predecessor and \(R\)-successor rules cannot close a cycle
through \(Q\), so both scans meet vertical boundaries. Between them
lies a finite strip with one start head and all blank payloads.
Vertical-bit propagation preserves this strip in time. The horizontal
boundary bit is uniform across it. The time shift is a permutation of
the finite cyclic index set, so the initial boundary row recurs. Before
its first return, every row is the exact deterministic machine update;
the head stays unique and cannot cross a boundary. The row before the
return permits a head only in state 215. Hence K215 halts on blank tape.

Conversely, any finite blank halting run has rectangular presentations
with independently padded space and time margins, followed by halt
padding. The standard helical presentation gives a cyclic word with
the four overlaps, one Start at the chosen origin, and no forbidden
\(E_*\). The detailed padding and actual-index tuning argument is
the package-06 blank-helical proof, whose hypotheses depend only on
the center predicate and these overlap identities. Removing corners
does not alter those hypotheses.

\section{Four exact, smaller coefficient recipes}
\label{sec:coefficients}
\subsection{A closed finite definition}

The following construction fixes all four integers in \eqref{eq:ports}.
There is no free program index, input, or oracle. Its finite local
predicate is the literal K215 predicate of Section~\ref{sec:local};
the complete transition table and predicate are included in the archive.
The coefficient definitions terminate after bounded finite enumeration,
although their binary expansions are infeasibly large.

Set \(a_0=439\) and \(\kappa=38\,256\,607\,897\).
Let \(W_0=S\), \(W_1=E_*\), with
\[
 S=(438,434,435,436,1),\qquad
 E_*=(438,0,437,438,437).
\]
List the remaining valid crosses in lexicographic order on
\(\{0,\ldots,438\}^5\), excluding these first two.
Coordinates are always prior, left, center, right, following.
The count \(\kappa\) is the exact count before the old compiler's
free-corner factor \(439^4\). The source census and the independent
replay are included; no halting test is performed in this enumeration.

Let \(\ell\) be the least integer at least two satisfying
\(2^\ell>\kappa+1\), and put \(A_0=2^\ell\). Define
\begin{align}
 p_0&=\max\left\{0,
       \left\lceil\frac{\kappa-2\ell-9a_0-4}{2}\right\rceil\right\},
 &j_{\max}&=5a_0+4+p_0,\notag\\
 \mu&=A_0-2+\sum_{j=1}^{j_{\max}}A_0^j,
 &m_0&=2\ell+14a_0+8+2p_0,\notag\\
 n_0&=m_0+1,
 &n_{\rm dum}&=n_0-\kappa-5a_0-4,\notag\\
 e_*&=\kappa+12a_0,
 &M_0&=\kappa+15a_0+n_{\rm dum}.
 \label{eq:cross-layout}
\end{align}
In this instance \(n_{\rm dum}=1\). The definitions retain the
original geometric separation between selectors, copies, dummies,
and anchors.

Use the ordered shape
\[
 \Sigma=((0,1),(1,0),(1,1),(1,2),(2,1)).
\]
For \((r,c)=\Sigma_z\), \(0\le z\le4\), and \(0\le u<a_0\), set
\begin{equation}\label{eq:cross-copy-address}
 e(r,c,u)=\kappa+3a_0+(2-r)3a_0+(2-c)a_0+u,
 \qquad j(z,u)=1+za_0+u.
\end{equation}
The allowed radix positions form the disjoint set
\begin{align}
 \mathcal E={}&\{0,\ldots,\kappa-1\}
 \ \cup\ \{e(r,c,u):(r,c)\in\Sigma,\ 0\le u<a_0\}\notag\\
 &{}\cup\{e_*,\ldots,e_*+n_{\rm dum}-1\}
 \ \cup\{M_0,3M_0,9M_0,27M_0\}.
 \label{eq:cross-support}
\end{align}
It has cardinality \(n_0\). The five copy blocks, in the order
\(\Sigma\), start at
\(\kappa+10a_0,\kappa+8a_0,\kappa+7a_0,\kappa+6a_0,\kappa+4a_0\);
each has length \(a_0\). All nonanchor positions lie below \(M_0\).

Define clause coefficients \(c_e\), for \(e\in\mathcal E\), by
\begin{align}
 c_i&=1+\sum_{z=0}^4 A_0^{\,1+za_0+W_i[z]}
            +\sum_{t=1}^4 A_0^{5a_0+t}
       &&(0\le i<\kappa),\notag\\
 c_{e(r,c,u)}&=A_0^{\,j(z,u)}
       &&((r,c)=\Sigma_z),\notag\\
 c_{3^{t-1}M_0}&=A_0^{5a_0+t}
       &&(1\le t\le4),\notag\\
 c_e&=0&&\text{at dummy positions}.
 \label{eq:cross-clause-coefficients}
\end{align}
The leading \(1\) counts selected crosses; the other digits pair each
selected cross with its five copies and four anchors.

\subsection{Band separation, radix, and the four ports}

Put
\begin{align}
 E_{\max}&=27M_0,&H_C&=51M_0+3a_0+1,\notag\\
 T_1&=H_C+2E_{\max}+a_0+1,
 &T_2&=T_1+4E_{\max}+2.
 \label{eq:cross-bands}
\end{align}
Let \(L_0\) be the least power of five at least
\(T_2+2E_{\max}+2\). For \(v>0\), let
\(\operatorname{bits}(v)=\lfloor\log_2v\rfloor+1\), and set
\begin{align}
 U_0&=4(n_0+2)\bigl(28(\kappa+5a_0+4)+6\bigr)+8,\notag\\
 b_{\min}&=\max\bigl\{
       \ell(j_{\max}+1)+3,\
       \ell(5a_0+4)+\operatorname{bits}(U_0)+1,\
       4\bigr\}.
 \label{eq:cross-radix-bound}
\end{align}
Let \(b\) be the least power of five at least \(b_{\min}\), and put
\[
 \rho=2^b,\qquad d=bL_0,\qquad B=\rho^{L_0}=2^d.
\]
All exponents used below are nonnegative.

The transport coefficients are
\begin{align}
 D_C={}&\rho^{3a_0}+\rho^{H_C+a_0}
       +\rho^{8M_0}+\rho^{24M_0}
       +\sum_{e\in\mathcal E}c_e
              \bigl(\rho^{T_1-e}+\rho^{T_2-e}\bigr),\notag\\
 D_R={}&\rho^{H_C}.
 \label{eq:cross-transport-constants}
\end{align}
Define the unshifted field mask by
\begin{align}
 M_F^0={}&\mu(\rho^{T_1}+\rho^{T_2})
       +\sum_{u=0}^{a_0-1}
            \bigl(\rho^{e(0,1,u)}+\rho^{e(1,1,u)}\bigr)\notag\\
 &+\sum_{u=0}^{a_0-1}
            \bigl(\rho^{H_C+e(1,0,u)}+\rho^{H_C+e(1,1,u)}\bigr)
       +\rho^{9M_0}+\rho^{27M_0}.
 \label{eq:cross-field-mask}
\end{align}
Finally set
\begin{equation}\label{eq:cross-word-mask}
 M_C=B-1-\sum_{\substack{e\in\mathcal E\\e\ne1}}\rho^e
                  -2\rho^{e_*},
 \qquad M_F=M_F^0+4.
\end{equation}
The four fixed polynomial coefficients are exactly
\begin{equation}\boxed{
 \beta=B-1,\qquad K=D_C+BD_R,\qquad M_C,\qquad
 \MF=M_F+B-1.}\label{eq:new-four-ports}
\end{equation}
In particular, the source port \(\MF\) is not the native mask \(M_F\).
Every arithmetic presentation in this article uses this same quartet.

The concrete scalar geometry is
\begin{center}\small
\begin{tabular}{@{}lr@{\qquad}lr@{}}
\toprule
Quantity&Value&Quantity&Value\\
\midrule
\(\ell\)&36&\(n_{\rm dum}\)&1\\
\(p_0\)&19\,128\,301\,935&\(j_{\max}\)&19\,128\,304\,134\\
\(m_0\)&38\,256\,610\,096&\(M_0\)&38\,256\,614\,483\\
\(e_*\)&38\,256\,613\,165&\(E_{\max}\)&1\,032\,928\,591\,041\\
\(H_C\)&1\,951\,087\,339\,951&\(T_1\)&4\,016\,944\,522\,473\\
\(T_2\)&8\,148\,658\,886\,639&\(b_{\min}\)&688\,618\,948\,863\\
\bottomrule
\end{tabular}
\end{center}
Thus
\begin{equation}\label{eq:smaller-d}
 b=5^{17},\qquad L_0=5^{19},\qquad
 d=5^{36}=14\,551\,915\,228\,366\,851\,806\,640\,625.
\end{equation}
The coefficient \(\beta\) has exactly \(d\) binary digits.
The corresponding package-06 value was \(5^{67}\); the reduction
factor is \(5^{31}\). This improves coefficient height but does not
make expansion feasible. The new recipe also removes four copy
coordinates and eight families of overlap tests.

\subsection{Populations, carry bounds, and exact convolution tests}

\begin{lemma}\label{lem:cross-interface}
The coefficients above satisfy
\begin{gather*}
 0<M_C,M_F<B-1,\qquad M_C\equiv2\pmod4,\qquad
 M_F\equiv4\pmod8,\\
 \pc(M_C)+\pc(M_F)=d.
\end{gather*}
They implement the four cross overlaps of
Section~\ref{sec:local} with the same carry-free native interface
as the package-06 compiler.
\end{lemma}

\begin{proof}
The powers of \(A_0\) occupy disjoint \(\ell\)-bit blocks, so
\[
 \pc(\mu)=\ell+5a_0+3+p_0,\qquad
 m_0=2\pc(\mu)+4a_0+2.
\]
The two clause sites in \(M_F^0\) contribute \(2\pc(\mu)\);
the four copy-test families contribute \(4a_0\); and the anchor
sites contribute two. Hence \(\pc(M_F^0)=m_0\) and
\(\pc(M_F)=m_0+1\). The added low bit \(4\) is disjoint from
all other support.

In \(M_C\), exactly \(n_0\) distinct bits are removed from \(B-1\):
one low bit at every native position except selector \(1\), plus
one upper bit at the first dummy. No borrow occurs. Thus
\[
 \pc(M_C)=d-m_0-1,\qquad
 \pc(M_C)+\pc(M_F)=d.
\]
Selector zero is permitted, selector one is forbidden, ordinary
native digits are Boolean, and the first dummy digit may be
\(0,1,2,3\). The lowest mask bits give the displayed congruences;
unused high positions give the strict range bounds.

The clause coefficients obey the conservative estimate
\[
 \sum_{e\in\mathcal E}c_e
       \le14(\kappa+5a_0+4)A_0^{5a_0+4},
 \qquad \mu<A_0^{j_{\max}+1}.
\]
Consequently the radix definition implies
\begin{equation}\label{eq:cross-carry}
 \rho\ge
 \max\left\{4(n_0+2)\left(2\sum_ec_e+6\right)+8,\
             2\mu+4,\ 16\right\}.
\end{equation}
This bounds every unshifted native field coefficient by
\(\rho/4-2\), including the permitted upper dummy bit.
Moreover \(T_2+E_{\max}<L_0\), so no unshifted field term
spills across a cell boundary.

At each clause site \(T_j\), convolution by \(D_C\) reads exactly
\(\sum_ec_ex_e\) from the current cell. The other clause band cannot
reach it because \(T_2-T_1>2E_{\max}\). The short-shift terms lie
strictly below the clause bands.

At \(e(0,1,u)\) and \(e(1,1,u)\), the short shift \(3a_0\)
reads, respectively, the current center and following copies.
The identity contribution from the temporal neighbor reads its
prior and center copies. These are the two temporal overlaps.

At \(H_C+e(1,0,u)\) and \(H_C+e(1,1,u)\), the shift \(H_C+a_0\)
reads the current center and right copies, whereas \(D_R\)
reads the spatial neighbor's left and center copies. These are
the two spatial overlaps.

At \(9M_0\) and \(27M_0\), the shifts \(8M_0\) and \(24M_0\)
read current anchors \(M_0\) and \(3M_0\); the temporal neighbor
contributes its anchors \(9M_0\) and \(27M_0\).
There are no unintended terms. Low support is below \(M_0\);
the four anchors are separated by powers of three; high clause
bands lie above the test sites; and
\(H_C=51M_0+3a_0+1\) separates horizontal tests from the possible
anchor sum \(24M_0+27M_0=51M_0\).
These support statements prove the exact convolution roles for
every word, not merely for the toy audits.
\end{proof}

\subsection{Why arbitrary bit rotations still synchronize}

We give the synchronization argument because it is the principal
compiler obligation beyond the population identity.
After arithmetic soundness recovers \(X=2^R\), write its binary
rotation exponent as \(bt+\epsilon\), \(0\le\epsilon<b\).
The possible rotated radix support is contained in a translate,
modulo \(L_0\), of
\[
 \mathcal E\cup\{e_*+1\}\subset[0,E_{\max}+1].
\]
The extra position records the upper dummy bit spilling into the
next radix position. Both cyclic distances between \(T_1\) and
\(T_2\) exceed \(E_{\max}+1\), so this rotated support cannot
contaminate both clause sites. One clause site is clean.

For \(\epsilon\le b-2\), a rotated digit is at most \(3\rho/4\).
For \(\epsilon=b-1\), it is at most \(\rho/2+1\).
Together with \eqref{eq:cross-carry}, neither case creates a carry.
Thus the mask tests apply to the actual field. More explicitly,
the transport equation and \(X=wq\equiv w\pmod{q-1}\) identify
\(F\) modulo \(q-1\) with the native field \(\widehat F\).
The carry bounds give \(0<\widehat F<q-1\), while arithmetic
soundness gives \(0<F<q\); therefore \(F=\widehat F\).

At the clean clause site, the low base-\(A_0\) digit counts
selected crosses. This count is at most \(\kappa<A_0-1\), and
its AND with \(A_0-2\) vanishes only for counts zero and one.
Every higher clause digit is a sum of the relevant selected
indicator and one copy or anchor bit. These sums are below \(A_0\),
so no base-\(A_0\) carry occurs. Their masked low bits force
evenness. Hence an occupied cell has exactly one selector, all
five matching copies, and all four anchors; an unoccupied cell
has none of these copies or anchors. Dummy bits have coefficient
zero and remain independent.

The marked origin is occupied. Its two anchor tests require
two distinct rotated digits to be odd. If
\(1\le\epsilon\le b-2\), all rotated digits are even.
If \(\epsilon=b-1\), only the single upper-dummy spill position
can be odd. Thus \(\epsilon=0\).

All nonanchor native positions lie below \(M_0\), and the only
pair of native positions separated by \(18M_0\) is
\((9M_0,27M_0)\). Since \(L_0>54M_0\), reduction modulo \(L_0\)
cannot supply another such pair through wraparound.
The two anchor tests therefore force \(t=0\pmod{L_0}\):
the surviving rotation is by a whole number of cells.

The copy tests now propagate occupancy to both spatial neighbors.
Indeed, a selected cross has one-hot copies, and its tested
center or side copy cannot match an empty neighbor. The spatial
offset \(1\) generates the whole cyclic word, so Start forces
every cell occupied. The four decoded overlaps then yield
\eqref{eq:reconstruct}. In particular a Start-containing solution
cannot have zero temporal stride, since its prior cell and center
are different. No uniqueness of Start, or independent End marker,
is required for soundness.

\subsection{Completeness at the actual packed index}

The preceding support argument also preserves the canonical
completeness construction of package 06~\cite{blank06}. We record
the index-control step explicitly, rather than assuming that a
formal congruent exponent can replace the actual packed value.

Choose a padded halting tableau with spatial width \(h_0\) and
temporal height \(H_0\), both sufficiently large powers of five,
such that \(H_0\ge25\) and \(N=h_0H_0>25d\).
The helical quotient has one Start at the origin. Set \(q=B^N\).
Initially set every dummy bit to zero and let \(C_0,F_0,R_0\)
be the actual native word, actual field, and actual packed index.
Switching on the low dummy at an eligible cell \(v\) changes
\begin{align}
 \delta C&=\rho^{e_*}B^v,\notag\\
 \delta F&=(D_C+BD_R+B^{h_0})\rho^{e_*}B^v,\notag\\
 \delta R&=-\Gamma B^v,\qquad
 \Gamma=\rho^{e_*}(q^2-1)
             [1+q(D_C+BD_R+B^{h_0})].
 \label{eq:cross-dummy-derivative}
\end{align}
Use distinct positions
\[
 \{4j:0\le j<N/5\}\cup\{1\}.
\]
Neither spatial nor temporal movement wraps these individual
contributions, because \(h_0=N/H_0\le N/25\).

The two clause-band terms in \(D_C\) cancel modulo five:
\[
 T_2-T_1\equiv2\pmod4,\qquad \rho\equiv2\pmod5.
\]
Since \(a_0\equiv3\pmod4\),
\begin{equation}\label{eq:cross-unit}
 2D_C-D_R
 \equiv2(2+3\rho^{H_C}+1+1)-\rho^{H_C}
 \equiv3\pmod5.
\end{equation}
Also \(q,B,B^{h_0}\equiv2\pmod5\), so the bracket in
\eqref{eq:cross-dummy-derivative} is \(2D_C-D_R\) modulo five.
Thus \(\Gamma\) is invertible modulo the power of five \(dN\).
No conditional correction monomial is needed.

Every residue modulo \(dN\) is a Boolean subset sum of the displayed
weights \(B^v\). To see this, \(B^4\) has order \(N/5\) modulo
\(dN\), by
\[
 \nu_5(2^{4dj}-1)=\nu_5(d)+1+\nu_5(j).
\]
Its powers enumerate exactly the residues \(1+5dt\), with
\(t\) modulo \(N/5\). For a desired residue, optionally subtract
the distinct weight \(B\) so that the remaining residue is
nonzero modulo five. Let \(k_0\) be its least positive residue
modulo \(5d\). Then \(1\le k_0<5d<N/5\) and \(k_0\) is a unit
modulo \(N/5\). A cyclic interval of \(k_0\) distinct residues \(t\)
has sum \(k_0t_0+k_0(k_0-1)/2\); varying \(t_0\) gives every
sum modulo \(N/5\). Choosing the corresponding powers of \(B^4\)
proves the subset assertion.

Apply it to
\((R_0-dh_0)\Gamma^{-1}\pmod{dN}\).
After this Boolean adjustment the \emph{actual} packed index obeys
\begin{equation}\label{eq:cross-tuned-index}
 R\equiv dh_0\pmod{dN}.
\end{equation}
The dummy is absent from every clause and copy test; its upper bit
remains zero. All masks survive the adjustment. The canonical
digit bounds also give
\[
 2C+F\le(\rho/4+1)\frac{q-1}{\rho-1}<q/3,
\]
so the outer slack \(q-F-C\) stays strictly positive.

The shifted field mask deserves a separate check. Its unit radix
digit comes only from the temporal neighbor's Start selector,
since \(D_C,D_R\) have positive minimum exponent. Canonical Start
occurs only at zero and \(0<h_0<N\). That digit is therefore zero
at the origin and zero or one elsewhere. The origin's two low
mask bits and the other cells' bit \(2\) vanish as required.
Dummy adjustment occurs at positive radix exponent and changes
none of these tests.

Consequently the inverse-population identity and the native masks
give, at this same final \(R\),
\[
 3q+1<R<q^4,\qquad
 R\equiv3\pmod4,\qquad
 \pc(R)=3dN+2.
\]
The positive half-binomial converse~\cite{halfbinomial,blank06}
now supplies fresh Pell witnesses at the actual index.
It gives \(X=2^R\), \(q^3\mid Y\), and hence \(w=X/q\in\Pos\).
Equation \eqref{eq:cross-tuned-index} yields
\(w\equiv B^{h_0}\pmod{q-1}\). The native transport quotient
\[
 z=\frac{(K+w)C-F}{q-1}
\]
is integral and positive, since \(w>q>F\) and \(C\ge1\).
This is the shifted transport convention of
\eqref{eq:system60}. The ordinary strong witness is obtained
from the canonical normalized one by multiplication by \(\Delta\).
Thus the smaller coefficient quartet preserves completeness.

\subsection{Scope and reproducible evidence}

The inherited dependencies are the marked helical-tableau argument,
the native carry and mask interpretation, and the half-binomial
Pell converse~\cite{blank06,raw75,halfbinomial}. The archive also
preserves the native 76-operation compiler proof, the half-binomial
mask correction, and the five-adic subset-control note, with exact
source identities. The present change is the five-coordinate
language and its revised finite layout; it does not silently
specialize a theorem whose input-marker hypotheses remain unmet.

The generator \code{compiler/cross_constants.py} emits the closed
recipe, all scalar geometry, and its audits. With \code{--tm}, it
checks the exact source-table hash, all thirteen census categories,
32 helical control presentations, and 1,144 projected crosses.
Its other checks include 26,106 cyclic lift checks, 100,384 exhaustive
small cross-word candidates, and 16,224 sparse convolution basis
tests. An independent audit rederives the large-instance geometry
and support isolation by interval arithmetic without importing the
generator.

These finite checks supplement the parametric proofs above.
They do not enumerate the actual selector alphabet, expand the
coefficients, exhibit a halting K215 run, or kernel-check the
Diophantine equivalence. All definitions and proofs are ordinary
finite mathematics formalizable in ZFC; supplying such an argument
is distinct from supplying an encoded ZFC derivation or a new
Lean/Rocq theorem.



\section{Unrestricted integer equations and exact degrees}
\label{sec:integer}

The fifteen-variable arithmetic proof required positivity of exactly
the nine coordinates in \eqref{eq:guarded}. Replace each of these by
its own disjoint copy of
\begin{equation}\label{eq:pos}
 \operatorname{Pos}(u,v,t)=1+u^2+v^2+t(t+1).
\end{equation}
This is always positive on integers and represents every positive
integer. For surjectivity, apply the three-square theorem to
\(4(n-1)+1\), which is not of the excluded form \(4^a(8b+7)\).
Exactly one of its three squares is odd; writing it as \((2t+1)^2\)
and the others as \((2u)^2,(2v)^2\) gives \eqref{eq:pos}.
The remaining six coordinates stay unrestricted. There are therefore
\(9\cdot3+6=33\) integer unknowns.

The quadratic system likewise requires positivity of only
\(J,F,\alpha,z,i,s,w,\eta,\zeta\). Its other 23 coordinates may
be arbitrary integers: the scale equations force the required
positive scales, the signed auxiliary lemma handles \(o,j,V,f,y\),
and the main-root argument handles \(D,\gamma\). Thus
\eqref{eq:pos} gives \(9\cdot3+23=50\) unrestricted unknowns.
Its residual degrees are at most 4 and its sum of squares has exact
degree 8. The first residual has a nonzero quartic leading form,
and a sum of real squares of leading forms cannot cancel identically.

\begin{center}\small
\begin{tabular}{@{}lrrrr@{}}
\toprule
Integer presentation & Unknowns & Equations & Degree & Operations\\
\midrule
Gap root norm&32&1&434&138\\
Root outer norm&32&1&680&136\\
Lifted \(\mathcal S_{60}\)&33&6&50 maximum&123\\
Lifted gap system&33&6&44 maximum&124\\
Lifted \(P_{69}\)&33&1&190&132\\
Lifted gap product&33&1&179&133\\
Lifted \(S_{77}\)&33&1&100&140\\
Lifted gap sum of squares&33&1&88&141\\
Lifted \(\mathcal Q_{32}\)&50&23&4 maximum&not optimized\\
Its sum of squares&50&1&8&not optimized\\
\bottomrule
\end{tabular}
\end{center}

\subsection{Degree certificates for the short circuits}

Give each positive witness degree one and fixed numerals degree zero.
Let \(Q=\beta J\), \(k_0=\eta+\zeta\), and
\(C_1=Q-F-\alpha\). The top parts of \(X,Y,a,c\) are
\(wQ,sQ^3,wsQ^4,k_0sQ^3\). The six factors of
\(P_{69}\) have exact degrees
\[
 22,\quad18,\quad28,\quad7,\quad2,\quad22.
\]
The main norm has its leading nominal terms cancel; expanding it gives
\[
 D^2-\Delta c^2
 =2ac(X+\gamma H)+(X+\gamma H)^2-Hc^2,
\]
whose leader is \(8\gamma (wsQ^4)^2(k_0sQ^3)\).
Multiplying the nonzero factor leaders gives
\begin{equation}\label{eq:Pleader}
 8Q^{55}h\gamma i^2k_0^9w^7s^{16}T^2f^4
       (wC_1-z_0Q),
\end{equation}
a nonzero degree-99 form for every admissible \(\beta>0\).
For the sum of squares, the auxiliary residual uniquely has degree 28,
so its squared leader proves exact degree 56.

Under the mixed integer conversion, the guarded variables have weight
two and the six signed variables weight one. The factor degrees become
\(44,35,50,13,4,44\), totaling 190; the maximum residual degree is 50.
Substitution of the disjoint leading forms \(u^2+v^2+t^2\) is injective:
set two coordinates in every triple to zero to obtain the injective
monomial map \(x\mapsto u^2\). Thus these weighted degrees are
attained, not merely upper bounds or random-specialization guesses.

\section{ZFC calibration and verification status}
\label{sec:calibration}

At the pinned snapshot, the accepted machine endpoint is
\begin{quote}\small
\code{KEndpoint215.blank_halts_iff_zfc_refutation}:\\
\code{BlankHalts machine215} \(\leftrightarrow\)\\
\code{BProv ZFCax_s [] Form.fBot}.
\end{quote}
The negated endpoint is
\code{KEndpoint215.nonhalting_iff_zfc_consistency}.
The verification dossier records Lean and Rocq builds for the same
literal 430-row machine; their source and table identities are distinct
from the new Diophantine work~\cite{k215,endpoint}.

The ordinary calibration argument turns this exact blank computation
into a finite refutation predicate for the specified ZFC presentation.
Standard effective proof presentations translate into one another by
finite proof translations formalizable in ZFC. Every construction above
uses finite combinatorics, integer arithmetic, induction, Pell sequences,
and explicitly bounded coefficient enumeration. The arithmetic and
compiler proofs, combined with the ordinary machine calibration,
are therefore formalizable in ZFC and give
\[
 \mathrm{ZFC}\vdash
 \bigl[\exists\text{ a solution}\bigr]
 \leftrightarrow
 \bigl[\exists p\;\operatorname{Proof}_{\mathrm{ZFC}}(p,\bot)\bigr].
\]
Negation gives \eqref{eq:target}. This does not infer ZFC-provability
from Lean's soundness; it uses the ordinary mathematical arguments
that the machine formalization mirrors.

No new kernel-checked proof of the native compiler, Pell bridge, or
present equivalence is included. The full Lean/Rocq builds were not
rerun. The computational checks establish exactly their advertised
finite properties: literal arithmetic identities, live gate counts,
degree witnesses, sparse convolution maps, small cyclic reconstruction,
the exact category census, and the source table hash. They do not
search for an accepting K215 run, materialize the giant coefficients,
or decide consistency.

No consistency assumption is needed to prove the biconditional. If
ZFC is consistent, its second incompleteness theorem prevents ZFC
from proving the resulting unsolvability statement. Mere consistency
alone should not be silently strengthened to a claim of two-sided
independence; an additional appropriate soundness hypothesis is needed
to exclude a proof of the false existential statement.

\section{Reproduction, source attribution, and further work}
\label{sec:reproduction}

The archive is self-contained for its finite checks and uses Python's
standard library. Run \code{python3 reproduce.py} from its root to
regenerate the arithmetic packets, quadratic residuals, coefficient
specification, and verification receipts. Run \code{make report} to
compile this article with pdfLaTeX. The main source and four included
TeX sections must stay together. A file manifest records SHA-256 values.
The fresh seven-stage run authenticates 24 exact upstream files and
emits 20 arithmetic presentations, including retained intermediate
tradeoffs. Its independent sparse-polynomial audit proves the six
auxiliary-gap identities by exact symbolic expansion, without sampling.
The independent interval audit covers every actual native position,
not just the small control instances.

The proof dependencies retain their original filenames and attribution.
This contribution uses the project's MIT No Attribution terms and
preserves the inherited K215 notices, including Andrew James Wade's
compiler ancestry. The principal mathematical credit is to the eight
source packages and the underlying native compiler development; the
present contribution is a checked synthesis with the additional
cross-language, transport, lift, and signed-domain simplifications
described explicitly above.

\subsection{Further research questions}

\begin{enumerate}
\item \textbf{Formalize one interface end to end.} A focused next theorem
is the signed auxiliary lemma with its exact hypotheses, followed by
the ordinary strong-rank lemma. The five-cell compiler can then target
the accepted K215 endpoint without importing any consistency premise.
\item \textbf{Quotient the prior-cell information.} The center predicate
reads the prior cell only through its horizontal-boundary bit. A
heterogeneous representation with four full-cell fields and two
boundary-bit copies would have only \(87\,923\,161\) local patterns.
The corresponding sparse layout and carry/synchronization proof have
not been completed here. This is a concrete next coefficient target,
not part of the claimed coefficient recipe.
\item \textbf{Directly compile the numerical calculus.} A local language
for the calibrated C222 proof search could avoid the 215-state machine
alphabet entirely. Such a change needs its own exact proof-predicate
calibration and cannot be justified only by an equiconsistency slogan.
\item \textbf{Optimize the fourteen-variable norm.} Its root resultant
has high degree. The ordinary strong norm offers smaller degrees but
does not permit the same immediate integral-square argument; any
attempt to use a rational discriminant must recover integrality.
\item \textbf{Reduce the remaining guarded coordinates.} The proof
already releases six signs in the fifteen-variable form. Releasing
\(i\) without excluding zero is invalid: it collapses the strong
and auxiliary conditions. A valid further reduction must address that
degenerate branch explicitly.
\item \textbf{Measure coefficient descriptions as well as gate counts.}
The five-cell change substantially reduces actual height while leaving
the gate count unchanged. Compact symbolic sums, effective smaller
alphabets, and bounds that charge coefficient bit length should be
reported alongside arithmetic operation totals.
\item \textbf{Internalize the equivalence in weaker arithmetic.} The
present claim is ZFC-provability by formalizable ordinary mathematics.
A claim in PA or a weaker system requires a separate arithmetization
and accounting for the induction used in the compiler and Pell lemmas.
\end{enumerate}

\begin{thebibliography}{99}\footnotesize\raggedright
\setlength{\itemsep}{1pt}\setlength{\parskip}{0pt}
\bibitem{collection} ProveIt contributors.
\emph{Con(ZFC) in Diophantine form: eight independent packages and
reconciliation}. Snapshot \pin.\\
\path{SetTheory/Cardinals/docs/reports/hilbert-tenth-problem/con-zfc-diophantine-forms/README.md}.
\bibitem{ordinary02} Package 02, \emph{A fixed Diophantine representation
of Con(ZFC)}, especially the ordinary strong norm and rank bootstrap.
\path{02-system63/report.tex} under the collection above.
\bibitem{root05} Package 05, \emph{Fixed-input Diophantine forms of ZFC
consistency}, root elimination and signed quotient proof.
\path{05-fixed15-16/paper/fixed_input_zfc.tex}.
\bibitem{blank06} Package 06, \emph{Explicit Diophantine equations for
Con(ZFC)}; \path{06-blank67/report.tex},
\path{proofs/fixed_blank_67_proof.md}, and
\path{proofs/blank_helical_proof.md}.
\bibitem{quartic08} Package 08, \emph{Explicit Diophantine equations for
the consistency of ZFC};
\path{08-start-only67/evidence/diophantine_zfc_encoding/start_only_quartic34_proof.md}.
\bibitem{rank} ProveIt contributors. \emph{Recovering integrality after
relaxing an auxiliary Pell norm}.\\
\path{Computability/HilbertTenthProblem/Papers/1980/PELL_RELAXED_AUXILIARY_PROOF.md}.
\bibitem{reverse} ProveIt contributors. \emph{Reversing the auxiliary
unit gives 89 operations at degree 160}.\\
\path{Computability/HilbertTenthProblem/Papers/research-wip/native-stream-queue/complete75_reversed_auxiliary89.md}.
\bibitem{asymmetric} ProveIt contributors. \emph{Asymmetric scale tradeoffs}.\\
\path{Computability/HilbertTenthProblem/Papers/research-wip/native-stream-queue/complete75_asymmetric_scale_tradeoffs.md}.
\bibitem{halfbinomial} ProveIt contributors. \emph{Half-binomial Pell kernel}.\\
\path{Computability/HilbertTenthProblem/Papers/research-wip/native-stream-queue/pell_kernel_half_binomial42.md}.
\bibitem{raw75} ProveIt contributors. \emph{A 75-operation fixed-index
universal certificate}.\\
\path{Computability/HilbertTenthProblem/Papers/1980/FIXED_RAW_UNIVERSAL_75_PROOF.md}.
\bibitem{k215} ProveIt contributors. \emph{Formal correctness of the exact
215-state machine}.\\
\path{Computability/BusyBeaver/Research/ZFC/k225-formal/verification/k215/README.md}.
\bibitem{endpoint} ProveIt contributors. \emph{K215 endpoint}.\\
\path{Computability/BusyBeaver/Lean/BusyBeaver/ZFC/KEndpoint215.lean}.
\bibitem{jones} J. P. Jones. Universal Diophantine equation.
\emph{Journal of Symbolic Logic} 47 (1982), 549--571.
\url{https://doi.org/10.2307/2273588}.
\bibitem{bayer2025} J. Bayer, M. David, M. Ha\ss ler, Y. Matiyasevich,
and D. Schleicher. \emph{Diophantine equations over \(\mathbb Z\):
universal bounds and parallel formalization}. arXiv:2506.20909v2, 2025.
\url{https://arxiv.org/abs/2506.20909v2}.
\end{thebibliography}
\end{document}
